\chapter{Infective Endocarditis}
\index{Infective Endocarditis}
\label{sec:Infective Endocarditis}

\section{Overview}


\begin{itemize}[noitemsep]
  \item Infective endocarditis is an infection of the endocardial surface of the heart, most commonly the heart valves
  \item Acute infective endocarditis is caused by virulent organisms (\textit{Staphylococcus aureus}) and progresses rapidly, while subacute disease is caused by less virulent organisms (viridans group streptococci).
\end{itemize}
\section{Causes}
This condition happens because of microbial infection of damaged or abnormal endocardium. Following exposure, the pathogen invades host tissues and replicates, triggering an immune response that contributes to the symptoms. The severity of infection depends on pathogen virulence, host immune status, and timely treatment. Cardiovascular disease develops through a complex interplay of genetic, environmental, and lifestyle factors. Atherosclerosis, characterized by plaque buildup in arterial walls, is a common underlying mechanism. Infection begins with pathogen entry through a susceptible portal (respiratory tract, gastrointestinal tract, skin, mucous membranes). The pathogen must overcome host defenses including physical barriers, innate immunity, and adaptive immune responses. Microbial virulence factors such as toxins, adhesins, and capsules enable the pathogen to establish infection and evade immune clearance. Host factors including age, nutritional status, immune competence, and comorbidities influence susceptibility and disease severity. The underlying mechanisms involve complex interactions between genetic predisposition and environmental factors. Disruption of normal physiological processes leads to the characteristic features of the condition. Multiple contributing factors may act synergistically to trigger or worsen the condition.
\section{Risk Factors}


\begin{itemize}[noitemsep]
  \item Close contact with infected individuals
  \item Compromised immune system
  \item Travel to endemic areas
  \item Poor sanitation and hygiene practices
  \item Lack of vaccination
  \item Exposure to contaminated food or water
  \item Advanced age
  \item Cigarette smoking
  \item Hypertension and diabetes
  \item Elevated cholesterol levels
  \item Family history of heart disease
  \item Obesity and physical inactivity
  \item Genetic predisposition and family history
  \item Advancing age
  \item Lifestyle factors (diet, exercise, smoking, alcohol)
  \item Environmental exposures
  \item Underlying medical conditions
  \item Medication use
\end{itemize}
\section{Signs and Symptoms}

\subsection{Common symptoms}
\begin{itemize}[noitemsep]
  \item You may experience fever, chills, night sweats, malaise, weight loss, new or changing cardiac murmur, splinter hemorrhages, Osler's nodes, Janeway lesions, Roth spots, and splenomegaly. Fever and chills Fatigue and malaise Localized pain and inflammation at infection site Swelling, redness, and warmth (local infection) Cough and respiratory symptoms (respiratory infections) Nausea, vomiting, and diarrhea (gastrointestinal infections) Headache and body aches Skin rash or lesions Enlarged lymph nodes Systemic symptoms in severe cases (hypotension, organ dysfunction)
  \item Pain or discomfort in the affected area
  \end{itemize}

\subsection{Emergency warning signs}
\begin{itemize}[noitemsep]
  \item Severe or worsening symptoms of infective endocarditis require immediate medical evaluation
\end{itemize}
\section{Progression and Stages}
The condition may progress through different stages of severity over time. Early stages often involve mild or subtle symptoms that may be overlooked. Moderate stages involve more noticeable symptoms affecting daily function. Advanced stages may involve significant impairment and complications.
\section{Complications}
\begin{itemize}[noitemsep]
  \item Disease progression leading to worsening symptoms
  \item Secondary infections or injuries
  \item Side effects of treatment
  \item Reduced quality of life
  \item Emotional and psychological impact
\end{itemize}
\section{Diagnosis}
To make the diagnosis, doctors look for modified Duke criteria: major criteria (positive blood cultures, evidence of endocardial involvement on echocardiography) and minor criteria (predisposing condition, fever, vascular phenomena, immunologic phenomena). Laboratory diagnosis includes pathogen identification through culture, PCR testing, or antigen detection. Serological tests can confirm exposure or immune response to the pathogen. Cardiovascular diagnosis relies on a combination of clinical history, physical examination, electrocardiography, imaging (echocardiography, CT angiography), and laboratory markers (troponin, BNP). Detailed history including exposure, travel, and vaccination status Physical examination focusing on the affected system Complete blood count with differential Microbiological culture of blood, sputum, urine, or wound PCR and nucleic acid amplification tests for rapid pathogen detection Antigen detection tests Serological testing for antibodies (IgM, IgG) Imaging studies (chest X-ray, CT, ultrasound) Gram stain and microscopy of clinical specimens Antimicrobial susceptibility testing to guide therapy

\begin{itemize}[noitemsep]
  \item Comprehensive medical history and physical examination
  \item Laboratory tests to assess organ function
  \item Imaging studies to visualize affected structures
  \item Specialized testing depending on the affected system
  \item Consultation with relevant specialists
\end{itemize}
\section{Prevention}
Hand hygiene with soap and water or alcohol-based sanitizer Vaccination according to recommended schedules Safe food handling and water purification Use of personal protective equipment in healthcare settings Safe sex practices including condom use Mosquito avoidance (nets, repellents, insecticides) Respiratory etiquette (covering coughs and sneezes) Isolation of infected individuals when indicated Prophylactic antibiotics or antivirals for high-risk exposures Travel precautions and pre-travel consultations

\begin{itemize}[noitemsep]
  \item Maintain a healthy lifestyle with balanced diet and exercise
  \item Avoid known risk factors
  \item Attend regular medical check-ups for early detection
  \item Follow recommended screening guidelines
\end{itemize}
\section{Treatment Options}
\begin{itemize}[noitemsep]
  \item Treatment focuses on prolonged intravenous antibiotics (4--6 weeks) based on culture and sensitivity
  \item Surgical intervention is indicated for heart failure, uncontrolled infection, large vegetations, or perivalvular abscess.
  \item Medications to manage symptoms and address underlying mechanisms
  \item Lifestyle modifications and self-care measures
  \item Physical therapy and rehabilitation
  \item Surgical intervention when indicated
  \item Regular monitoring and follow-up care
\end{itemize}
\section{Living with It}
\begin{itemize}[noitemsep]
  \item Follow your treatment plan as prescribed
  \item Maintain regular follow-up with your healthcare provider
  \item Make necessary lifestyle adjustments
  \item Seek support from family, friends, or support groups
  \item Stay informed about your condition
\end{itemize}
\section{Outlook}
Your outlook depends on the causative organism, patient factors, and timeliness of intervention. Most patients recover fully with appropriate treatment, though some infections may lead to long-term complications. Outcomes depend on the pathogen, host immune status, and timeliness of treatment. With proper management and lifestyle modifications, many patients maintain good quality of life. Long-term outcomes depend on the extent of disease, adherence to treatment, and control of risk factors. Most infectious diseases resolve completely with appropriate treatment, especially when diagnosed early. Outcomes depend on the pathogen, host immune status, and timeliness of intervention. Some infections can become chronic (HIV, hepatitis B/C, herpes) requiring long-term management. Antimicrobial resistance may complicate treatment and worsen outcomes. Public health measures and vaccination programs continue to reduce the global burden of infectious diseases.
